\subsection{Where compatible divisor corrections can be supported}

\begin{statement}{A.89}{Proved}
{The exact local ideal-norm kernels of the radical tower}
Retain A.86 and write $\mathcal D_{n,q}$ for the free
$\Z_\ell$-module on the primes of $K_n$ above a finite prime
$q$, with transition map the ideal norm.  Put
\[
 W_{n,q}=\ker\bigl(\mathcal D_{n,q}
                                  \xrightarrow{\mathrm N_{n,0}}\Z_\ell\bigr),
 \qquad W_{\infty,q}=\varprojlim_n W_{n,q}.
 \tag{A.89.1}
\]
For $q\ne\ell$ let $m_q=v_q(A)\in\Z$ be its ordinary
valuation.  If $m_q=0$, then
\[
                                  W_{\infty,q}=0.
 \tag{A.89.2}
\]
If $m_q\ne0$, put $s_q=v_\ell(m_q)$.
Choose compatible root labels so that tame inertia acts
on $\Z_\ell$ by translation by $m_q$ and Frobenius acts
by $t\mapsto qt+\alpha_q$.
Let $f_1,\ldots,f_{h_q}$ be the lengths of its cycles on
$\Z_\ell/m_q\Z_\ell$.  These are the actual residue degrees;
their prime-to-$\ell$ factors are retained.  There is a natural
identification, after choosing an ordering of these cycles,
\[
 W_{\infty,q}\simeq
 \ker\left(\Z_\ell^{h_q}\longrightarrow\Z_\ell,
                      \quad(t_i)_i\longmapsto\sum_i f_i t_i\right).
 \tag{A.89.3}
\]
The transition maps on the prime basis are identities once
$n\ge s_q$: each subsequent extension is totally ramified
at these primes, so its relative residue degree is one.

In the dyadic range $A/3\in(\widehat{\Q^\times})^2$ of A.77,
one has $W_{\infty,2}=0$ for $\ell=2$.
For $\ell=3$ and $v_3(A)=1$, one has $W_{\infty,3}=0$.
The inverse kernel on the real sign coordinates in A.86 is
also zero when $\ell=2$.

There is an injective actual coordinate map
\[
 \varprojlim_n W_n\longrightarrow
      \prod_{q<\infty}W_{\infty,q}\ \times W_{\infty,\infty}.
 \tag{A.89.4}
\]
Its image consists exactly of the families with the following
condition: for every fixed $n$ and $j$, all but finitely many
rational primes $q$ have zero level-$n$ vector modulo $\ell^j$.
This condition is part of the assertion; an arbitrary product
family need not lie in the source of (A.89.4).
Consequently, for $\ell=2,3$ in the stated range, every compatible
divisor correction in (A.86.5) is supported on the primes
$q\ne\ell$ with $v_q(A)\ne0$, through the explicit modules
(A.89.3) and this convergence condition.
\end{statement}
\begin{proof}
The root labels and the prime modules come from the actual
finite Kummer torsors, so we give their transition maps explicitly.
Fix a compatible system of $\ell$-power roots of unity and
the generators $\alpha_n$ of A.86.  Label the roots of
$T^{\ell^n}-a_n$ by $t\in\Z/\ell^n\Z$ using
$\zeta_{\ell^n}^{\,t}\alpha_n$.
Their maps to the preceding torsor are the maps in (A.86.1);
on the labels they are reduction modulo $\ell^n$.
The Kummer cocycle of $A$ gives the affine action
$t\mapsto\chi_\ell(g)t+a(g)$ on the inverse limit
$\Z_\ell$.
At $q\ne\ell$ its tame inertia is translation by $m_q$;
wild inertia and the other tame factors have trivial image.
The cyclotomic character is unramified, and an arithmetic
Frobenius gives $t\mapsto qt+\alpha_q$.
These formulas also show that changing compatible labels
conjugates the affine action and leaves the cycle description
unchanged.

At level $n$, decomposition-group orbits are the primes of
$K_n$ above $q$.  Within one such orbit, inertia orbits have
size the ramification index, and Frobenius permutes those
inertia orbits in a cycle of length equal to the residue degree.
If a prime at level $n+1$ has residue degree $f'$ over
$\Q_q$ and lies over one of residue degree $f$, its ideal
norm is multiplication by the integer $f'/f$ on that prime.
Thus the prime-module transition is the actual ideal norm,
not a division by a cycle length.

\emph{Unramified affine orbits.}
Suppose first $m_q=0$.
The root tower is unramified at $q$, and its primes are the
cycles of $\varphi(t)=qt+\alpha_q$ on $\Z/\ell^n\Z$.
Every finite orbit of $\varphi$ on $\Z_\ell$ is fixed.
Indeed, if $\varphi^h(t)=t$ for $h>0$, then
\[
 0=(1+q+\cdots+q^{h-1})\bigl((q-1)t+\alpha_q\bigr).
 \tag{A.89.5}
\]
The first factor is a nonzero ordinary integer, and $\Z_\ell$
is an integral domain.  Hence $(q-1)t+\alpha_q=0$.
Since $q-1\ne0$, there is at most one fixed point
$t_0\in\Z_\ell$.

We need the integral norm compatibility, not just this set of
finite orbits.  Fix a level-$n$ cycle $O$ whose inverse image
in $\Z_\ell$ does not contain $t_0$; if no fixed point exists,
this means any cycle.  The $\ell$-adic valuations of the
cycle lengths above $O$ tend uniformly to infinity with the
level.  Otherwise there would be a constant $c$ and arbitrarily
large levels $r$ with a point above $O$ of period having
$\ell$-part at most $\ell^c$.
The prime-to-$\ell$ part of the order of any affine permutation
of $\Z/\ell^r\Z$ divides $\ell-1$, since that affine group has
order $(\ell-1)\ell^{2r-1}$.
Thus these points satisfy
$\varphi^{(\ell-1)\ell^c}(t)=t\pmod{\ell^r}$.
Compactness of the inverse image of $O$ gives a limiting
point satisfying this equality in $\Z_\ell$, contrary to
(A.89.5) and the choice of $O$.

For a compatible integral divisor, its coefficient at $O$
is the sum over all cycles $O'$ above it of
\[
             \frac{|O'|}{|O|}\,t_{O'},\qquad t_{O'}\in\Z_\ell.
 \tag{A.89.6}
\]
Every summand is divisible by an arbitrarily high power of
$\ell$ once the upper level is large enough.  Hence that
coefficient is zero.  If the fixed point exists, the only
remaining coefficient at each level is on its singleton
cycle.  These coefficients agree under the transition maps,
so the inverse divisor group is one copy of $\Z_\ell$,
generated by that degree-one branch.  Its norm to $\Q_q$
is the identity.  If the fixed point does not exist, the
inverse divisor group is zero.  In either case the inverse
norm kernel is zero, proving (A.89.2).

\emph{Nonzero ordinary valuation.}
Now $m_q\ne0$ and $s=s_q<\infty$.
For $n\ge s$, the inertia orbits on $\Z/\ell^n\Z$ are
the fibres of reduction to $\Z_\ell/m_q\Z_\ell$.
Each has size $\ell^{n-s}$.
The affine Frobenius permutes this one fixed finite set;
write its cycle lengths as $f_1,\ldots,f_{h_q}$.
The corresponding primes have
\[
             e(\mathfrak q_{n,i}/q)=\ell^{n-s},\qquad
             f(\mathfrak q_{n,i}/q)=f_i,\qquad
                         \sum_i f_i=\ell^s.
 \tag{A.89.7}
\]
At every subsequent level exactly one prime lies above
each of them, its ramification index is multiplied by
$\ell$, and its residue degree is unchanged.
The ideal norm on this stable prime basis is therefore
the identity.  Its norm to the rational-prime divisor is
the weighted sum $\sum_i f_i t_i$.  This proves (A.89.3),
including the prime-to-$\ell$ factors of the weights.

\emph{The two coefficient places and real signs.}
For $\ell=3$, the representatives in A.77 have
$v_3(a_n)=1$.  The polynomial $T^{3^n}-a_n$ is Eisenstein
over $\Q_3$.  Thus there is one prime, with residue degree
one, and every ideal norm in the tower is the identity on
its $\Z_3$ coefficient.  The kernel is zero.

For $\ell=2$, use the full local decomposition group proved
in the dyadic part of A.77.  The local Galois closure has
group $(\Z/2^n\Z)\rtimes(\Z/2^n\Z)^\times$, and the root
field is the fixed field of the stabilizer
$(\Z/2^n\Z)^\times$.
An index-two subgroup containing that stabilizer is the
kernel of a nonzero homomorphism to $\Z/2\Z$ trivial on
the stabilizer.  On the cyclic translation group there is
exactly one such homomorphism, reduction modulo two.
Thus the root field has exactly one quadratic subfield,
namely $\Q_2(\sqrt{a_n})=\Q_2(\sqrt3)$, which is ramified.
Its degree is $2^n$.  If its residue degree were greater
than one, its maximal unramified subfield would contain
an unramified quadratic subfield, a contradiction.
Consequently the unique prime has residue degree one.
Every subsequent relative residue degree is one, proving
the claimed zero local inverse norm kernel in the full scope.

For the constructed ghosts one also has $v_2(A)=0$,
and in that special case there is a direct polynomial
check.  Remove a $2^n$th power of two from $a_n$ locally;
the remaining unit has squareclass three and is
$3\pmod8$.  Replacing $T$ by $T+1$ makes the constant
coefficient have valuation one and all intermediate
binomial coefficients even.  The resulting Eisenstein
polynomial gives the same residue-degree assertion.

The same dyadic assumption makes every rational $a_n$
positive.  Choose positive real $\alpha_n$, so the rational
transition factors $b_n$ in (A.86.1) are positive.
At each level $n\ge1$ the two real embeddings send
$\alpha_n$ to its positive and negative real roots.
Only the positive embedding has real extensions at the
next level; it has two, and their signs multiply under
the real norm.  On additive sign vectors the transition is
\[
                (u_+,u_-)\longmapsto(u_++u_-,0)
                       \quad\text{in }(\Z/2\Z)^2.
 \tag{A.89.8}
\]
A compatible family therefore has $u_-=0$ at every level,
and a constant $u_+$, whose norm to the sign over $\Q$
is itself.  The inverse kernel on signs is zero as claimed.

\emph{The actual global restricted family.}
At every level, a completed divisor is specified by its
vectors at the rational primes, subject to (A.86.6), and
by its finite sign vector when required.
The norm to $\Q$ preserves each rational-prime component
and the sign component separately.  Therefore a member
of $W_n$ has its $q$ component in $W_{n,q}$ and its sign
component in the corresponding sign kernel.
This gives (A.89.4), and the map is injective because
these coordinates determine every level-$n$ divisor.

For a family in its target, the stated convergence
condition is exactly the condition that its level-$n$
prime coefficients belong to (A.86.6).
There are at most $[K_n:\Q]$ primes above each rational
prime, so using rational primes in the convergence
condition is equivalent to using all primes of $K_n$.
Assemble the coordinates at each level.  The result lies
in $W_n$, and the local norm compatibilities give the
global norm compatibility.  This proves the asserted
description of the image, without replacing it by an
unrestricted product.  Equations (A.89.2), (A.89.3), and
the coefficient-place and sign calculations prove the
final support assertion.
\end{proof}

In the constructed ordinary ghosts, this theorem identifies
the precise local modules from which a correction of the
compatible ideal-class family in A.86 can be assembled.
It proves neither that their class images generate that
family nor that the full one-variable obstruction vanishes.

\noindent The exact finite certificate
\path{code/norm_local_divisor_certificate.py}, with retained output
\path{code/norm_local_divisor_certificate_output.txt}, checks the
stable inertia quotients and all residue-degree weights, sample
uniform affine-orbit norm growth, the shifted Eisenstein equations,
and the real sign transition.  The global convergence condition
is proved above and is retained in the statement.
